\pdfbookmark[1]{Abstract}{Abstract}
\chapter*{Abstract}

A number theoretic transform is correct only if its modulus is prime and its twiddle root has order
exactly the transform length. Both facts are proved here at the point of use, by one Proth witness
for primality and two exponentiations for the order, on unbounded integers at any width. A
1{,}233-digit modulus is proved prime in under three tenths of a second, with the witness recorded
beside the verdict. A root of half the required order passes every invariant computable from the
twiddle table in \(O(n)\), and the order test is the only one that separates it. Two silent wrong
answers fall to the same discipline: a composite modulus of the correct Proth shape, and a 32-bit
overflow in a device butterfly.

An integer relation search is measured as a function of its precision. Its cost grows with an
exponent rising to \(2.64\) per doubling, close to the \(2.585\) a Karatsuba multiply on operands
linear in the precision predicts, while the lattice reduction grows no faster than the multiply
count. The precision at which a device multiply overtakes the host is one at which a single search
cell takes about a day, and no faster device closes that gap.

The record of configurations states each claim with the result that would refute it, written before
the verdict, and the constraints every entry is bounded by.
\cleardoublepage
